\section{Results}
\label{sec:results}

\subsection{RQ1: VLM Baseline Performance on Misviz}
\label{sec:rq1}

Table~\ref{tab:misviz_main} summarizes binary classification performance across all four baseline conditions. Claude Haiku~4.5 substantially outperforms Qwen3-VL-8B on both binary F1 (0.830 vs.\ 0.723 in the vision-only condition) and per-type F1 (PT-F1). Both models achieve reasonably high recall but lower precision, indicating a tendency to over-flag charts. Importantly, high binary F1 masks weak per-type accuracy: the Claude vision-only baseline achieves binary F1~=~0.830 but PT-F1~=~39.3\% (computed over types with non-zero support), and flags 42\% of non-misleading charts as misleading.

\begin{table}[h]
  \centering
  \caption{Misviz benchmark --- binary classification results (271 charts). PT-F1 is the macro-average F1 over misleader types with non-zero support.}
  \label{tab:misviz_main}
  \small
  \begin{tabular}{llccccc}
    \toprule
    Model & Condition & Accuracy & Precision & Recall & Binary F1 & PT-F1 \\
    \midrule
    Claude Haiku 4.5 & vision\_only & 0.771 & 0.782 & 0.883 & 0.830 & 39.3\% \\
    Claude Haiku 4.5 & vision\_text & 0.779 & 0.832 & 0.813 & 0.822 & ---    \\
    Qwen3-VL-8B      & vision\_only & 0.672 & 0.773 & 0.678 & 0.723 & ---    \\
    Qwen3-VL-8B      & vision\_text & 0.664 & 0.845 & 0.573 & 0.683 & ---    \\
    \bottomrule
  \end{tabular}
\end{table}

\paragraph{Effect of text context.}
Adding textual grounding had no statistically significant effect on Claude (F1 drops slightly from 0.830 to 0.822) and substantially hurt Qwen's recall (0.678 $\to$ 0.573). This is a negative result: contrary to the multi-modal fusion intuition, injecting axis label and value text into the VLM prompt does not improve, and often degrades, misleader detection. We attribute this to two effects: (1) text context increases the input length and appears to dilute the model's attention to visual features, and (2) OCR extraction errors introduce false evidence that the model treats as authoritative.

\paragraph{Per-type analysis.}
Table~\ref{tab:pertype} shows per-type F1 for the two Claude conditions. Both conditions excel on structurally salient types---dual axis (F1 0.769/0.732) and 3D distortion (0.696/0.698)---while failing on the ``blind-spot'' types that require quantitative axis analysis: inconsistent tick intervals (0.0/0.111), discretized continuous variable (0.0/0.0), and inconsistent binning (0.095/0.118). Inappropriate axis range (0.154/0.111) and inappropriate item order (0.222/0.0) are also weak. This confirms a systematic gap between visually-obvious and quantitatively-grounded misleader types.

\begin{table}[h]
  \centering
  \caption{Per-type F1 for Claude Haiku~4.5 on Misviz (271 charts). Types are grouped by detection difficulty.}
  \label{tab:pertype}
  \small
  \begin{tabular}{lcc}
    \toprule
    Misleader Type & Vision-only & Vision+text \\
    \midrule
    \multicolumn{3}{l}{\textit{Visually salient types}} \\
    Dual axis              & 0.769 & 0.732 \\
    3D distortion          & 0.696 & 0.698 \\
    Pie chart misuse       & 0.571 & 0.526 \\
    Line chart misuse      & 0.424 & 0.529 \\
    Truncated axis         & 0.388 & 0.500 \\
    Inverted axis          & 0.348 & 0.300 \\
    Misrepresentation      & 0.329 & 0.308 \\
    \midrule
    \multicolumn{3}{l}{\textit{Quantitative / blind-spot types}} \\
    Inappropriate item order    & 0.222 & 0.000 \\
    Inappropriate axis range    & 0.154 & 0.111 \\
    Inconsistent binning        & 0.095 & 0.118 \\
    Inconsistent tick intervals & 0.000 & 0.111 \\
    Discretized continuous var. & 0.000 & 0.000 \\
    \midrule
    \textbf{Macro PT-F1}   & \textbf{0.333} & \textbf{0.328} \\
    \bottomrule
  \end{tabular}
\end{table}

\subsection{RQ2: Effect of Tool Augmentation and Pipeline Design}
\label{sec:rq2}

Table~\ref{tab:pipeline} reports PT-F1, precision, recall, and false positive/negative counts for all pipeline variants evaluated on the 271-chart Misviz set (Claude Haiku~4.5).

\begin{table}[h]
  \centering
  \caption{Pipeline variant comparison on Misviz benchmark (271 charts, Claude Haiku~4.5). PT-F1 is macro-averaged per-type F1. VLM calls counts the average API calls per chart.}
  \label{tab:pipeline}
  \small
  \begin{tabular}{lcccccc}
    \toprule
    Architecture & PT-F1 & PT-Prec & PT-Rec & FP & FN & VLM calls \\
    \midrule
    VLM-only baseline           & 39.3\% & 35.7\% & 43.7\% & 157 & 112 & 1 \\
    V3 (OCR+Rules injected)     & 38.9\% & 31.3\% & 51.3\% & 224 &  97 & 1 \\
    V4 (baseline + CLEAN veto)  & 43.5\% & 44.3\% & 42.7\% & 107 & 114 & 1 \\
    V5 (V4 + DePlot rules)      & 37.6\% & 30.0\% & 50.3\% & 233 &  99 & 1 \\
    V6 (targeted re-ask)        & 43.0\% & 36.5\% & 52.3\% & 181 &  95 & 2 \\
    V7 (sequential re-ask)      & 40.1\% & 30.1\% & 59.8\% & 276 &  80 & ${\sim}7$ \\
    V8 (V7 + 3-vote majority)   & 38.9\% & 29.1\% & 58.8\% & 285 &  82 & ${\sim}19$ \\
    \textbf{V4 + ViT veto (ours)} & \textbf{45.2\%} & \textbf{50.0\%} & \textbf{41.2\%} & \textbf{82} & 117 & \textbf{1} \\
    \bottomrule
  \end{tabular}
\end{table}

\paragraph{Tool injection hurts (V3, V5).}
V3 injected OCR-derived rule verdicts directly into the VLM prompt. PT-F1 dropped from 39.3\% to 38.9\% and false positives increased by 43\% (157~$\to$~224). The OCR ``flagged'' verdicts had only 0--4\% accuracy, introducing misleading evidence. V5 replicated this failure with DePlot: extracted data tables generated 20 false positives per 7 true positives from rule application, driving PT-F1 down to 37.6\%.

\paragraph{Post-processing veto improves results (V4).}
V4 applied OCR rule verdicts as a post-hoc veto: only when rules confidently classify a type as ``clean'' (99--100\% precision) do we override the VLM's positive prediction. Combined with disambiguation examples for the three most-confused types, this reduced false positives to 107 and raised PT-F1 to 43.5\%---the first meaningful gain over baseline.

\paragraph{Multi-call re-asking reveals attention limitation (V6, V7).}
V7's per-type sequential calls achieved the highest recall (59.8\%) and fewest false negatives (80), confirming that VLMs \textit{can} detect blind-spot types when asked in isolation. However, the ${\sim}7\times$ call overhead also inflated false positives to 276, yielding lower PT-F1 (40.1\%) than V4.

\paragraph{Self-consistency fails (V8).}
Majority-voting across three V7 runs (temperature~0.7) degraded PT-F1 further to 38.9\%, confirming that chart misleader errors are \textit{systematic} rather than random---sampling cannot correct them.

\paragraph{ViT selective veto achieves best PT-F1 at baseline cost.}
The final pipeline adds a ViT-B/16 selective veto to V4, removing false positives on five high-precision types (tick intervals, binning, axis range, item order, truncated axis) without additional VLM calls. This reduces false positives from 107 to 82 and raises PT-F1 to \textbf{45.2\%} with precision \textbf{50.0\%}---a 5.9-point improvement over the VLM-only baseline at identical inference cost.

Table~\ref{tab:pertype_improvement} shows per-type F1 gains from the final pipeline over baseline.

\begin{table}[h]
  \centering
  \caption{Per-type F1 change from VLM-only baseline to V4\,+\,ViT veto (selected types with notable changes).}
  \label{tab:pertype_improvement}
  \small
  \begin{tabular}{lccc}
    \toprule
    Misleader Type & Baseline F1 & V4\,+\,ViT F1 & $\Delta$ \\
    \midrule
    Inconsistent binning        & 9.5\%  & 34.5\% & $+$25.0\% \\
    Inconsistent tick intervals & 0.0\%  & 20.5\% & $+$20.5\% \\
    Truncated axis              & 38.8\% & 53.1\% & $+$14.3\% \\
    3D distortion               & 69.6\% & 80.8\% & $+$11.2\% \\
    Misrepresentation           & 32.9\% & 41.7\% & $+$ 8.8\% \\
    \bottomrule
  \end{tabular}
\end{table}

\subsection{RQ3: Generalization to Real SEC Filings}
\label{sec:rq3}

Table~\ref{tab:sec_main} reports binary detection performance on the 96-chart SEC validation set. All 10 companies in our dataset had active SEC comment letters, so every flagged chart constitutes a true positive (precision~=~1.0 for all conditions). The key differentiator is recall: how many charts from companies with known violations does the model flag?

\begin{table}[h]
  \centering
  \caption{SEC 10-K validation set --- binary detection results (96 charts, 10 companies). All companies have active SEC comment letters (all ground-truth positive), so precision = 1.0 for all conditions.}
  \label{tab:sec_main}
  \small
  \begin{tabular}{llccc}
    \toprule
    Model & Condition & Recall & F1 & Charts Flagged \\
    \midrule
    Claude Haiku 4.5 & vision\_only & 0.375 & 0.545 & 36/96 \\
    Claude Haiku 4.5 & vision\_text & 0.396 & 0.567 & 38/96 \\
    Qwen3-VL-8B      & vision\_only & 0.167 & 0.286 & 16/96 \\
    Qwen3-VL-8B      & vision\_text & 0.073 & 0.136 & 7/96  \\
    \bottomrule
  \end{tabular}
\end{table}

\paragraph{Cross-domain gap.}
The best condition (Claude vision+text) achieves F1~=~0.567 and recall~=~0.396 on the SEC set, compared to binary F1~=~0.822 on Misviz. The gap reflects two distributional differences: SEC financial charts are predominantly simple bar, line, or stock performance charts with fewer extreme visual distortions than the Misviz benchmark; and our ground truth (the presence of an SEC comment letter) is a coarser signal than the Misviz per-chart annotations.

\paragraph{Company-level variation.}
Table~\ref{tab:sec_ticker} shows per-company flag rates for Claude vision-only. Performance is highly uneven. STZ and UPS show very high flag rates (100\% and 80\%), consistent with their comment letters flagging numerous Non-GAAP presentation issues. CHPT shows 0\% flagging despite having a comment letter; manual inspection confirms its charts are predominantly simple trend lines without obvious structural misleaders, suggesting the comment letter concerns were textual rather than visual. BCO (50\%), VERI (57\%), and IRBT (20\%) fall in between.

\begin{table}[h]
  \centering
  \caption{Per-company chart flag rates for Claude Haiku~4.5 (vision-only) on SEC 10-K validation set. All companies have SEC comment letters (has\_gt\_violation = true).}
  \label{tab:sec_ticker}
  \small
  \begin{tabular}{lcc}
    \toprule
    Ticker & Charts & Flag Rate \\
    \midrule
    STZ   & 10 & 100\% \\
    UPS   & 10 &  80\% \\
    VERI  &  7 &  57\% \\
    BCO   & 10 &  50\% \\
    KMI   &  9 &  22\% \\
    IRBT  & 10 &  20\% \\
    LITE  & 10 &  20\% \\
    LYFT  & 10 &  20\% \\
    OGN   & 10 &  10\% \\
    CHPT  & 10 &   0\% \\
    \bottomrule
  \end{tabular}
\end{table}

\paragraph{Qualitative validation.}
Beyond binary detection, the SEC pipeline performs Non-GAAP compliance analysis on extracted financial tables. For the companies where the comment letter specifically cited Non-GAAP presentation issues (STZ, UPS, VERI), the pipeline flagged the same tables and correctly identified the relevant regulatory provisions (Regulation~G and Item~10(e) of Regulation~S-K). This qualitative agreement with SEC comment letter content provides partial validation that the system's reasoning is grounded in the correct compliance framework, not merely pattern-matching on visual features.

\subsection{Summary}
\label{sec:summary}

Table~\ref{tab:summary} consolidates the key findings across all three research questions.

\begin{table}[h]
  \centering
  \caption{Key findings summary.}
  \label{tab:summary}
  \small
  \begin{tabular}{p{3cm}p{9cm}}
    \toprule
    Research Question & Finding \\
    \midrule
    RQ1: VLM capability & Claude Haiku~4.5 achieves binary F1~=~0.830 but PT-F1~=~39.3\%, revealing a large gap between binary and fine-grained detection. Blind-spot types (tick intervals, binning, discretized variable) score near zero. \\
    \midrule
    RQ2: Tool augmentation & Injecting OCR and rule outputs into VLM prompts consistently degrades PT-F1. Using the same tools as post-processing vetoes (V4) improves PT-F1 to 43.5\%. Adding a ViT-B/16 selective veto achieves the best result of 45.2\% PT-F1 at baseline inference cost (1 VLM call/chart). \\
    \midrule
    RQ3: SEC generalization & The best pipeline achieves recall~=~0.396 and F1~=~0.567 on 96 real SEC charts, with large per-company variance. STZ and UPS are flagged at 80--100\%; companies with only textual violations (CHPT) are missed entirely, reflecting a genuine gap in visual-vs.-textual misleader detection. \\
    \bottomrule
  \end{tabular}
\end{table}
